% beamerposter hard-codes hmargin=1cm regardless of paper size (its line 176),
% which on A0 is a 1.2 % margin — far too tight, and the reason `size=a3'
% cannot reproduce these proportions. Re-issue \geometry AFTER it loads.
\geometry{paperwidth=841mm,paperheight=1189mm,
          top=28mm,bottom=28mm,head=0pt,headsep=0pt,foot=0pt}

% Re-issuing \geometry does NOT override beamerposter's hard-coded hmargin=1cm
% — the horizontal margin has to be set through beamer itself. 10 mm is far
% too close to the trim edge for a printed A0; many large-format printers
% cannot even reach it.
\setbeamersize{text margin left=22mm,text margin right=22mm}

\newlength{\colgap}   \setlength{\colgap}{24mm}
% \colw is set inside the frame, not here: in the preamble \textwidth still
% holds beamerposter's value and the columns come out 10 % too wide.
\newlength{\colw}     \setlength{\colw}{240mm}
\newlength{\blockskip}\setlength{\blockskip}{15mm}
\newlength{\headskip} \setlength{\headskip}{22mm}

\setlength{\parskip}{0.5\baselineskip}
\setlength{\parindent}{0pt}

% Only complain about overflow that is actually visible.
\hfuzz=2pt

% The height a column may occupy: the sheet minus the top and bottom margins,
% minus the header band and the gap under it.
\newlength{\colmax}
\setlength{\colmax}{\dimexpr\paperheight-28mm-28mm-165mm-\headskip\relax}
\newsavebox{\colbox}
% With [t] alignment the reference point sits on the first line, so almost the
% whole box lives in its DEPTH — \ht alone reports zero. Total extent is
% height plus depth.
\newcommand{\showcol}[1]{%
  \typeout{COLHEIGHT: column #1 = \the\dimexpr\ht\colbox+\dp\colbox\relax
           \space of \the\colmax}}
